\documentclass[../../../include/open-logic-section]{subfiles}

\begin{document}

\olfileid{sth}{story}{urelements}
\olsection{సమితి కాని మూలకాలు ఉండాలా?}

తరువాతి కొన్ని అధ్యాయాల్లో
సమితుల సంచిత-దశలవారీ
భావన నుంచి స్వీకృతాలను
రాబట్టడానికి ప్రయత్నిస్తాం.
కానీ ముందుకు వెళ్లేముందు
\emph{సమితి కాని మూలకాల}
(urelements) గురించి
మరింత చెప్పాలి.

\olref[approach]{sec}లోని
చిత్రం పాత \emph{సమితుల}
నుంచి మాత్రమే కొత్త
సమితులను ఏర్పరచడానికి
అనుమతించింది. అయితే
కొన్ని \emph{సమితులు కాని}
వస్తువులు---ఆవులు,
పందులు, ఇసుక రేణువులు,
లేదా మరేవైనా---సమితుల
\tetoken{మూలకాలుగా}{!!{element}s}
ఉండవచ్చని అనుమతించాలనుకోవచ్చు.
అప్పుడు కొన్ని ప్రాథమిక
మూలకాలతో, అంటే \emph{సమితి
కాని మూలకాలతో}, ప్రారంభించి,
ప్రతి దశ $S$లో ఆ మూలకాలు
లేదా $S$కు ముందు దశల్లో
ఏర్పడిన సమితులు ఏ కలయికలోనైనా
ఉన్న ``సాధ్యమైన అన్ని''
సమితులను ఏర్పరుస్తామని
చెబుతాం. అప్పుడు చిత్రం
ఇలా ఉంటుంది:
\begin{center}
	\begin{tikzpicture}[scale=0.6]
	\tikzset{cut_here/.style={densely dotted}}
	\draw (-4,6) -- (-1, 0) -- (1,0) -- (4,6); 
	\draw[cut_here] (-5,8)--(-4,6);
	\draw[cut_here] (5,8)--(4,6);
	\draw(-1.5, 1)--(1.5, 1);
	\draw(-2, 2)--(2, 2);
	\draw(-2.5, 3)--(2.5,3);
	\draw(-3, 4)--(3, 4);
	\draw(-3.5, 5)--(3.5, 5);
	\draw(-4, 6)--(4, 6);
	\node[label] at (2, 0) {\small 0};
	\node[label] at (2.5, 1) {\small 1};
	\node[label] at (3, 2) {\small 2};
	\node[label] at (3.5, 3) {\small 3};
	\node[label] at (4, 4) {\small 4};
	\node[label] at (4.5, 5) {\small 5};
	\node[label] at (5, 6) {\small 6};
	\end{tikzpicture}
\end{center}
ఇప్పుడు మనం ఒక నిర్ణయం
తీసుకోవాలి: \emph{సమితి కాని
మూలకాలను అనుమతించాలా?}

తాత్త్వికంగా చూస్తే,
సమితి కాని మూలకాలను
మన సిద్ధాంతంలో చేర్చడం
సమంజసం. ప్రధాన కారణం
సమితి సిద్ధాంతాన్ని
\emph{అన్వయించగలిగేలా}
చేయడమే. దీనిని చూపడానికి,
రెండు సమితులు $A$,~$B$
సమాన పరిమాణం గలవని,
అంటే $\cardeq{A}{B}$ అని,
వాటి మధ్య ద్విగుణ ప్రమేయం
ఉన్నప్పుడే చెబుతామని
\olref[sfr][siz][]{chap} నుంచి
గుర్తుచేసుకోండి. ఇప్పుడు
పొలంలోని ఆవులూ, దొడ్డిలోని
పందులూ సమితులు ఏర్పరుస్తే,
``ఆవులెన్నో పందులూ అన్నే''
అనే వాదనను సమితి సిద్ధాంతంతో
వివరించవచ్చు. కానీ సమితి కాని
మూలకాలను నిషేధించి,
ఆవులూ పందులూ \emph{సమితులు
ఏర్పరచవు} అంటే, అలాంటి
సమితి-సిద్ధాంత వివరణ
అందుబాటులో ఉండదు.
నిజానికి సమితులు కాకుండా
ఇతర దేనికైనా సమితి సిద్ధాంతాన్ని
నేరుగా అన్వయించడం కష్టమవుతుంది.
(సమితి కాని మూలకాలను
చేర్చడానికి మరిన్ని కారణాల కోసం
\citealt[pp.~vi, 24,
50--1]{Potter2004} చూడండి.)

గణితంలో అయితే సమితి కాని
మూలకాలను అనుమతించడం
అరుదు. కొంతవరకు, అవి
లేకుండా సమితి సిద్ధాంతాన్ని
రూపొందించడం \emph{చాలా స్వల్పంగా}
సులభం కావడమే కారణం.
కానీ కొన్నిసార్లు సమితి కాని
మూలకాలను సమితి సిద్ధాంతం
నుంచి తొలగించడానికి
ఆసక్తికరమైన సమర్థనలూ
కనిపిస్తాయి:
\begin{quote}
	సమితి సిద్ధాంతం గణితానికి
	పునాది అనే నమ్మకం ప్రకారం,
	సమితుల గురించే మాట్లాడుతూ
	గణితం మొత్తాన్ని పట్టుకోగలగాలి.
	అందువల్ల మన చరం ఆవులు,
	పందుల వంటి వస్తువులపై
	వ్యాపించకూడదు.
	%కానీ $C$ ఒక ఆవు అయితే, $\{C\}$ ఒక సమితి; అయినా అది సముచిత గణిత వస్తువు కాదు.
	\citep[p.~8]{Kunen1980}
\end{quote}
అందువల్ల అన్వయయోగ్యతపై
దృష్టి సమితి కాని మూలకాలను
\emph{చేర్చాలని} సూచిస్తుంది;
గణితాన్ని శుద్ధ సమితి
సిద్ధాంతానికి తగ్గించే
పునాదివాద లక్ష్యంపై దృష్టి
వాటిని \emph{తొలగించాలని}
సూచించవచ్చు. స్వల్ప సోమరితనం
కూడా వాటిని తొలగించే
వైపే నడిపిస్తుంది.

మనం అత్యంత సులభమైన
దారినే అనుసరిస్తాం.
అయితే కొంతవరకు దీనికి
బోధనాపర సమర్థన ఉంది.
సమితి సిద్ధాంత తత్త్వశాస్త్రాన్ని
చదవడం ప్రారంభించడానికి
అవసరమైన సమితి సిద్ధాంత
మూలాంశాలను పరిచయం చేయడమే
మన లక్ష్యం. ఆ తాత్త్విక
సాహిత్యంలో ఎక్కువ భాగం
సమితి కాని మూలకాలు
\emph{లేకుండా} రూపొందించిన
సమితి సిద్ధాంతాలనే
చర్చిస్తుంది. కాబట్టి ఆ
సాహిత్యానికి దగ్గరగా ఉంటే
ఈ పుస్తకం మరింత
ఉపయోగకరంగా ఉండవచ్చు.

\end{document}
